\section{The original ordinary Hecke algebra and its residual block}
\label{sec:ordinary-hecke-algebra-v042}

\status{This chapter proves the original \texttt{hiComm} and
\texttt{hiAlg}, using the actual ordinary modules and operators.
The base coefficient ring may have nilpotents. The domain hypothesis
is used only for the source's torsion-freeness conclusion. No
arithmetic-point separation or automorphy theorem is inferred.}

Let \(H\) be the original \texttt{HSet}, with its original
\texttt{Ax}, coefficient DVR \(O=c.O\), and an \(O\)-algebra \(A\)
of class \texttt{IsCO}. This means that \(A\) is a complete
Noetherian local ring, the map \(O\to A\) is local, and it induces
an isomorphism on residue fields. In particular the residue field
of \(A\) is finite, so \(A\) is of class \texttt{IsCNL}.
Fix the original character \(\theta\) satisfying \texttt{IsWt},
and fix any original auxiliary level \(K\). Put
\[
 U=H.\mathrm{lev}(K),\qquad
 Y=H.X.\mathrm{Ord}(H.\tau_A(A),\theta,U).
\]
OST19 makes \(Y\) finite free. The operator index set is literally
\[
       H.\mathrm{Op}=(H.W'\times\{1,\ldots,n\})
                    \sqcup(H.V_\ell\times\{1,\ldots,n\}).
\]
The first kind of operator is the integral spherical double-coset
sum at \(w\), acting through \(H.\mathrm{sec}(w)\) and the original
away action. The second is \(U_{a_{v,j}}\), where \(a_{v,j}\) has
\(\operatorname{diag}(\varpi_v I_j,I_{n-j})\) at \(v\) and identity
elsewhere. In the source the displayed \(j\) is represented by
an element of \(\operatorname{Fin}(n)\) plus one.

For a prescribed function \(a:H.\mathrm{Op}\to O\), write
\(a_i^A\) for its image in \(A\), and define the actual residual
part
\begin{equation}\label{eq:oha-locpart}
 L=\left\{y\in Y:
   \ \forall i\ \forall k\ge0\ \exists r\ge0,\
               (T_i-a_i^A)^r y\in\mathfrak m_A^kY\right\}.
\end{equation}
This is the source \(H.MK(\theta,a,K)\). Its algebra \(H.TK\)
is the \(A\)-subalgebra of \(\operatorname{End}_A(L)\) generated
by the restrictions of these same \(T_i=H.\mathrm{op}(A,i)\).
The powers in \eqref{eq:oha-locpart} are allowed to depend on
\(i,k,y\); the proof below does not replace this condition by
semisimplicity.

% BEGIN OHA1
\paragraph{OHA1. Integral finite double-coset sums on arbitrary representations.}
\label{claim:OHA1}
Let \(G\) be a group, \(C\subset G\) a subgroup, and assume
\(CgC/C\) is finite for every \(g\). For any representation
\(\rho:G\to\operatorname{Aut}_A(M)\) and \(v\in M^C\), set
\[
                 [CgC]v=\sum_{xC\in CgC/C}\rho(x)v .
\]
This is independent of the representatives: replacing \(x\)
by \(xc\) fixes its summand. Left multiplication by \(C\)
permutes the summands, so the result belongs to \(M^C\).
An operator commuting with every \(\rho(x)\) commutes with
this sum.

More generally, a finitely supported bi-\(C\)-invariant integer
function on double cosets acts by
\(\sum_{xC}f(x)\rho(x)\). Expanding two finite sums and collecting
the coefficient of each \(zC\) gives
\[
 (f*h)(z)=\sum_{xC\in G/C}f(x)h(x^{-1}z).
\]
Hence this action respects convolution. Each coefficient is an
integer, and the same identity holds over every commutative \(A\)
by the map \(\mathbb Z\to A\); no averaging denominator is used.
These sums are exactly the original \texttt{dc} on fixed vectors.

If \(G\) is a topological group and \(C\) is compact and open,
the finiteness assumption holds: \(G/C\) is discrete, and the
continuous image of the compact set \(C\) under \(c\mapsto cgC\)
is finite. This argument does not require the representation
\(\rho\) to have any topology.
% END OHA1

% BEGIN OHA2
\paragraph{OHA2. Spherical commutativity over every coefficient ring.}
\label{claim:OHA2}
Let \(F\) be one of the original nonarchimedean local fields and
\(C=\operatorname{GL}_n(\mathcal O_F)\). All its spherical
double-coset sums commute on \(C\)-fixed vectors, over every
commutative coefficient ring.

Here is a direct integral proof. Smith reduction over the DVR
gives \(g=u d v\) with \(u,v\in C\) and \(d\) diagonal. To obtain
this for a matrix with fractional entries, first clear their
finitely many denominators; the integral Smith form then gives
the asserted form after dividing its diagonal by that common
denominator. Consequently
\[
       g^t=v^t d u^t
          =(v^tu^{-1})g(v^{-1}u^t)\in CgC .
\]
Thus transpose preserves every double coset and reverses
multiplication.

For explicit measure control, on the open set
\(\operatorname{GL}_n(F)\subset M_n(F)\) the measure
\(|\det g|_F^{-n}\,dg\), with \(dg\) additive matrix Haar
measure, is both left and right invariant. Indeed multiplication
by \(h\) on the left or right has additive Jacobian
\(|\det h|_F^n\), canceled by the determinant factor.
Transpose preserves this measure. Normalize it by \(C\) having
volume one. For compactly supported bi-\(C\)-invariant integer
functions its convolution is the finite sum in OHA1, so
transpose induces the anti-automorphism
\[
                       (f*h)^t=h^t*f^t .
\]
The previous double-coset calculation gives \(f^t=f\), \(h^t=h\)
and \((f*h)^t=f*h\). Hence \(f*h=h*f\), an identity of integer
functions. OHA1 transfers it to every \(A\)-representation.
In rank zero the group is trivial and the conclusion is immediate.

The diagonal with no uniformizer entries is the identity, and
the one with all \(n\) entries equal to the uniformizer is central.
Their double cosets each have one right coset. Thus the extra
source identities \(\mathrm{sphOp}_0=1\) and
\(\mathrm{sphOp}_n=\rho(\varpi I_n)\) follow as well.
This proves the entire \texttt{sphComm} interface.
% END OHA2

% BEGIN OHA3
\paragraph{OHA3. The original spherical operators preserve the actual level.}
\label{claim:OHA3}
For \(w\in H.W'\), put \(s_w=H.\mathrm{sec}(w)\) and
\(C_w=\operatorname{GL}_n(\mathcal O_{L_w})\).
The original \texttt{W\_le} gives \(s_w(C_w)\subset U\).
Thus a \(U\)-fixed function is \(C_w\)-fixed under the
representation \(g\mapsto r_a(s_w(g))\).

For \(u\in H.UQ\), \texttt{W\_dec} gives the literal factorization
\[
          u=s_w(c)u',\qquad c\in C_w,\quad u'\in H.UQ,
          \qquad [u',s_w(g)]=1\quad\text{for every }g .
\]
OHA1 makes the spherical image \(C_w\)-fixed. The action of
\(u'\) commutes term by term with its finite sum and fixes the
initial function. Hence \(u\) fixes the spherical image.
Every auxiliary-level subgroup also fixes that image, since
\texttt{sectionsCommuteAtSplit} gives the required commutation.
The stabilizer of a vector is a subgroup, so it contains their
join, exactly \(U=H.UQ\vee\bigvee_q H.\mathrm{sec}(w_q)(K_q)\).
This proves level preservation with the source order of factors.
No unproved normality of that join is used.
% END OHA3

% BEGIN OHA4
\paragraph{OHA4. Stability and pairwise commutation on the actual ordinary module.}
\label{claim:OHA4}
Every \(H.\mathrm{op}(A,i)\) preserves \(Y\), and all these
operators commute there.

For spherical stability, the source away action is
\[
       (r_a(u)f)(g)=\tau_A(u)\,f(g(u,1)).
\]
It preserves left rational invariance, upper-unipotent invariance,
the prescribed torus character and uniform lower-unipotent
continuity. For the last assertion its finitely many matrix
coefficients multiply differences already in the same ideal
\(\mathfrak m_A^k\). It also preserves every value filtration
and commutes with all \(p\)-side translations, hence with every
\(U_b\). For \(f\in Y\), OST10 supplies at every precision a
period \(r\ge1\) with \(U_z^rf-f\) in that value filtration.
Any finite sum of away translates satisfies the same period
condition. OHA3 gives the remaining level condition, proving
spherical stability in the original \(Y\).

Each \(a_{v,j}\) belongs to the positive monoid, by its decreasing
list of exponents and the original product construction.
OST7 and OST20 give its ordinary stability and invertibility.
Such operators commute with the spherical sums term by term.
Two positive operators commute by OST5 and commutativity of the
diagonal monoid.

For two spherical places, the exact \texttt{Wm\_mul\_Wm\_comm}
field has two alternatives. If their section images commute,
their two finite sums commute term by term. Otherwise it gives
an isomorphism \(e:\operatorname{GL}_n(L_{w'})
\simeq\operatorname{GL}_n(L_w)\) preserving the integral compact
groups, with \(s_{w'}=s_w\circ e\). It induces a bijection of
the relevant right-coset sets. OHA1 therefore identifies the
second sum, on fixed vectors, with the double-coset operator
at \(e(a_{w',j'})\) in the first group. OHA2 applies to arbitrary
double cosets, so these two sums commute. Applying OHA3 before
the second use of the coset identification checks its fixed-vector
hypothesis at both stages. This proves the full original
\(\mathrm{CommOn}(Y,H.\mathrm{op}(A))\).
% END OHA4

% BEGIN OHA5
\paragraph{OHA5. A finite commutative operator algebra over a complete local ring.}
\label{claim:OHA5}
Let \(A\) be a complete Noetherian local ring, let \(Y\) be a
finite \(A\)-module, and let \((T_i)_{i\in I}\) be any commuting
family in \(\operatorname{End}_A(Y)\). The index set may be
infinite. The algebra
\[
                      T=A[T_i:i\in I]\subset\operatorname{End}_A(Y)
\]
is finite over \(A\). If \(Y\ne0\), it has a finite product
decomposition
\[
                 T=\prod_{\mathfrak q}T_{\mathfrak q},
                 \qquad 1=\sum_{\mathfrak q}e_{\mathfrak q},
\]
where each factor is a complete Noetherian local ring, finite
over \(A\), and its maximal ideal lies over \(\mathfrak m_A\).
The factors are the localizations at all maximal ideals of \(T\).

Choose finitely many generators of \(Y\). Evaluation on them
injects \(\operatorname{End}_A(Y)\) into a finite power of \(Y\).
Noetherianity therefore makes this endomorphism module finite,
and its submodule \(T\) is finite as well. The commuting hypothesis
makes \(T\) commutative. Finiteness implies integrality, so every
maximal ideal of \(T\) contracts to \(\mathfrak m_A\).
Now apply the finite-algebra decomposition for a complete
Noetherian local base (Stacks Project, Lemma 10.160.2, Tag 0325).
Its hypotheses have just been checked on this exact \(A,T\).
Likewise every finite \(A\)-module used below is complete and
separated, by Lemma 10.97.1, Tag 00MA, since \(A=\widehat A\).

One can see explicitly why the projectors exist. The finite
dimensional algebra \(T/\mathfrak m_AT\) is a product of local
Artinian algebras. A lift \(b\) of any one of its idempotents
is successively corrected by
\[
                 b\longmapsto b-(b^2-b)(2b-1)^{-1}.
\]
The error is squared at every step and \(2b-1\) is a unit,
because its square is \(1+4(b^2-b)\). Completeness supplies
the lift, including in residue characteristic two. Lifts are
unique: if commuting idempotents \(e,f\) are congruent modulo
\(\mathfrak m_AT\), the idempotents \(e(1-f)\) and \(f(1-e)\)
are in the Jacobson radical and hence zero. Orthogonality and
the identity sum therefore lift. The resulting factors have
one maximal ideal each. This also explains the actual
idempotents in the displayed decomposition.
% END OHA5

% BEGIN OHA6
\paragraph{OHA6. The original adic definition selects exactly one idempotent block.}
\label{claim:OHA6}
In OHA5 fix any \(a_i\in A\) and use the definition of \(L\)
in \eqref{eq:oha-locpart}. There is at most one maximal ideal
\(\mathfrak q_a\) of \(T\) for which every \(T_i-a_i\) belongs
to \(\mathfrak q_a\). If it exists, then
\[
                            L=e_{\mathfrak q_a}Y .
                            \tag{OHA.1}
\]
If it does not exist, \(L=0\). Thus in both cases the projector
belongs to the original algebra \(T\).

If \(Y=0\), take the zero projector and all assertions are immediate.
Otherwise use the decomposition in OHA5. To prove uniqueness, consider
\(J=\mathfrak m_AT+(T_i-a_i:i\in I)\).
The quotient \(T/J\) is generated over \(A/\mathfrak m_A\)
by the already specified scalar images of its generators.
It is therefore either zero or the field \(A/\mathfrak m_A\).
In the nonzero case \(J\) itself is the unique maximal ideal
in question.

Decompose \(Y=\bigoplus_{\mathfrak q}Y_{\mathfrak q}\) with
\(Y_{\mathfrak q}=e_{\mathfrak q}Y\). In the selected factor,
the maximal ideal is nilpotent modulo \(\mathfrak m_AT_{\mathfrak q_a}\).
Consequently some exponent \(h\ge1\) satisfies
\(\mathfrak q_a^hT_{\mathfrak q_a}
 \subset\mathfrak m_AT_{\mathfrak q_a}\).
For every \(i,k\), \((T_i-a_i)^{hk}\) maps
\(Y_{\mathfrak q_a}\) into \(\mathfrak m_A^kY_{\mathfrak q_a}\).
This proves the inclusion from right to left in (OHA.1).

For any other factor choose \(i\) with \(T_i-a_i\notin\mathfrak q\).
This element is a unit in \(T_{\mathfrak q}\), and its inverse
preserves \(\mathfrak m_A^kY_{\mathfrak q}\).
If \(y\in L\), its \(\mathfrak q\)-component therefore belongs
to \(\mathfrak m_A^kY_{\mathfrak q}\) for every \(k\).
Separatedness of this finite \(A\)-module forces it to be zero.
There are only finitely many factors, giving (OHA.1).

This is also the literal source \texttt{locPart}: for fixed
\(i,k\), the inverse images of \(\mathfrak m_A^kY\) under
successive powers of \(T_i-a_i\) are increasing, because
that operator preserves \(\mathfrak m_A^kY\). Their submodule
supremum is consequently their union. Its elementwise statement
is exactly \eqref{eq:oha-locpart}.
Any endomorphism commuting with every \(T_i\) commutes with
the projector, since it commutes with every element of \(T\).
No choice of a semisimple generic fiber has entered this proof.
% END OHA6

% BEGIN OHA7
\paragraph{OHA7. All fields of the original residual-part freeness theorem.}
\label{claim:OHA7}
The residual part \(L\) is finite, stable under the original
operators, has a stable complement in \(Y\), and satisfies
\(\mathrm{locPart}(L,T,a)=L\). If \(Y\) is projective over \(A\),
then \(L\) is finite free over \(A\).

Use \(L=eY\) and \(Z=(1-e)Y\) from OHA6. They are complementary
finite modules and are stable because \(e\in T\). For every ideal
\(I\subset A\),
\[
                       L\cap IY=IL:
\]
if \(y\) is in the left side, applying \(e\) to a finite expression
for \(y\in IY\) gives \(y=ey\in IL\); the reverse inclusion is
immediate. Powers of \(T_i-a_i\) preserve \(L\), so this equality
at \(I=\mathfrak m_A^k\) changes its residual-part condition
relative to \(Y\) into that relative to \(L\). Hence localization
is idempotent.

A direct summand of a projective module is projective. To recall
the finite local freeness step, choose lifts of a basis of
\(L/\mathfrak m_AL\). Nakayama makes the resulting map
\(A^r\to L\) surjective. Projectivity splits it. The complementary
kernel has zero reduction modulo \(\mathfrak m_A\), and is finite,
so Nakayama makes it zero. This gives a basis of \(L\).
The pairwise commutation relations restrict from \(Y\).
These statements are exactly the five conclusions of
\texttt{LocFree.locFree}.
% END OHA7

% BEGIN OHA8
\paragraph{OHA8. The whole original \texttt{hiComm} endpoint.}
\label{claim:OHA8}
For every original \(K\), the original operators preserve and
commute on \(Y\), and \(Y\) is finite free. For every original
residual function \(a:H.\mathrm{Op}\to O\), they also preserve
and commute on \(H.MK(\theta,a,K)\), which is finite free.

Indeed \texttt{IsCO} and the finite residue field of the
coefficient DVR give \texttt{IsCNL} on this exact \(A\).
The \texttt{fin} field of \(H.\mathrm{Ax}\) and the
\texttt{cont/free} fields of \(H.\mathrm{IsWt}\) give the
hypotheses of OST19. OHA4 proves the remaining commutation
and stability. Apply OHA7 to the actual operators and the
images of the prescribed \(a_i\). This works when that block
is zero as well. It proves every quantifier of \texttt{hiComm},
not just the assertion for a selected residual eigensystem.
% END OHA8

% BEGIN OHA9
\paragraph{OHA9. The actual block algebra is complete local with the prescribed residue.}
\label{claim:OHA9}
In OHA5--OHA7 assume \(L\ne0\). Let \(T_L\subset
\operatorname{End}_A(L)\) be generated by the original restricted
operators. Then \(T_L\) is finite over \(A\), complete Noetherian
local, its residue field is exactly \(A/\mathfrak m_A\), and
the map \(A\to T_L\) is local. Each \(T_i-a_i\) belongs to
its maximal ideal. If multiplication by \(b\in A\) is injective
on \(Y\), then its scalar image is a non-zero-divisor on \(T_L\).

The selected factor in OHA6 exists because \(L\ne0\).
Restriction identifies \(eT=T_{\mathfrak q_a}\) with \(T_L\).
It is onto because its images are the original generators.
It is injective because an element \(et\) killing \(eY\)
also kills \((1-e)Y\), and \(T\subset\operatorname{End}_A(Y)\)
acts faithfully. All local and completeness claims now follow
from OHA5. The residue is \(T/\mathfrak q_a=A/\mathfrak m_A\)
by the proof of uniqueness in OHA6; this also proves the
prescribed generator congruences.

For the last assertion, if \(bt=0\) as an endomorphism of \(L\),
then \(b\,t(y)=0\) for every \(y\in L\subset Y\).
Injectivity on \(Y\) gives \(t(y)=0\) for all \(y\), hence \(t=0\).
This proves the entire \texttt{HeckeLocal.heckeLocal} conclusion,
including its version for any faithfully embedded algebra whose
image is \(T_L\): transport along that specified algebra
isomorphism, keeping its scalar map.
% END OHA9

% BEGIN OHA10
\paragraph{OHA10. Inverses already belong to the finite operator algebra.}
\label{claim:OHA10}
Let \(M\) be a finite module over any commutative ring \(A\),
and let \(u\) be an \(A\)-linear automorphism. Then
\(u^{-1}\in A[u]\). In particular every subalgebra containing
\(u\) contains its inverse.

Choose finitely many generators of \(M\) and express \(u^{-1}\)
on them by a matrix over \(A\). Its monic characteristic
polynomial annihilates \(u^{-1}\) on \(M\): this follows by
applying the adjugate identity to that matrix and then the
surjection from the free module on the chosen generators.
If \(M=0\) the conclusion is immediate. Otherwise write this
annihilating polynomial as
\[
                    X^r+c_{r-1}X^{r-1}+\cdots+c_0,\qquad r\ge1 .
\]
Multiplication of its identity at \(u^{-1}\) by \(u^{r-1}\)
gives
\[
            u^{-1}=-c_{r-1}-c_{r-2}u-\cdots-c_0u^{r-1}.
\]
This proof uses no freeness, reducedness or characteristic
assumption on \(A\).

Consequently, if \(L\subset Y\) is finite and stable under
an operator invertible on finite \(Y\), its inverse preserves
\(L\), by the same polynomial formula on \(Y\).
Its restriction is invertible on \(L\), and applying the formula
on \(L\) puts the restricted inverse in the algebra generated
by that restricted operator. This is the precise unit argument
needed for the original Hecke generators.
% END OHA10

% BEGIN OHA11
\paragraph{OHA11. The full original \texttt{hiAlg} endpoint.}
\label{claim:OHA11}
Assume \(H.MK(\theta,a,K)\ne0\). Its original algebra
\(H.TK(\theta,a,K)\) satisfies all six conclusions of
\texttt{HiAlgUnit.hiAlg}:
\begin{enumerate}
\item It is an \(O\)-algebra of class \texttt{IsCO}.
\item Its actual scalar map from \(A\) is finite.
\item If \(A\) is a domain, multiplication by any nonzero
  scalar from \(A\) is injective on this algebra.
\item Every generator differs from the original prescribed
  \(a_i\in O\) by an element of its maximal ideal.
\item Every original positive-side generator \(U_{a_{v,j}}\)
  is a unit in this algebra.
\item These very generators generate the entire algebra over \(A\).
\end{enumerate}

OHA8 supplies the finite free ordinary module and residual block.
OHA9 gives the finite complete local algebra, its local scalar
map and the congruences. The composite \(O\to A\to T_L\) is
local. The two residue maps are surjective, so every element
of \(T_L\) is congruent to the image of an element of \(O\).
Together with Noetherianity and completeness, these are
exactly the defining \texttt{IsCO} fields. If \(A\) is a domain,
its nonzero scalars act injectively on the free module \(Y\);
OHA9 gives the stated algebra torsion-freeness. Reducedness or
domain assumptions were not used for any earlier conclusion.

OST20 gives invertibility of every \(U_{a_{v,j}}\) on \(Y\).
OHA10 first restricts its inverse to \(L\), and then places that
inverse in \(T_L\). Hence the source \texttt{gen} itself is a
unit; this is stronger than invertibility after passing to a
fraction field. Finally \(T_L\) was defined as the algebra
generated by the restricted operators. The images of
\(\mathrm{gen}_i\) are precisely those restrictions, so the
subalgebra they generate internally is all of \(T_L\).
No field-valued eigensystem or arithmetic point was introduced.
% END OHA11

% BEGIN OHA12
\paragraph{OHA12. Installation on the same arithmetic Hecke setting at every level.}
\label{claim:OHA12}
For the original \(H=\mathrm{hset}(Q,\tau_0)\), FNL16 supplies
all its \texttt{Ax} fields and OTW12 supplies the same character
\(\theta\), fixed before all \(K\), with its full
\texttt{IsWt}. Therefore OHA8 and OHA11 apply to every original
level and every original residual function \(a\), whenever the
source nonzero-block hypothesis for OHA11 holds.

In particular the original \texttt{hiComm} and \texttt{hiAlg}
are now available after the original finite-level and ordinary
structure theorems. These statements do not require the
integral \texttt{TypeForAlt} producer. They do not supply
nonvanishing of a chosen block, an arithmetic specialization
map, a Galois determinant, dense arithmetic points, or the
four Hida outputs on one ideal. Those are subsequent producer
obligations. The coefficient embedding, level, Hecke
normalizations and character have remained the same throughout.
% END OHA12

\subsection{Sources and scope of theorem reuse}

The original interfaces are \texttt{HiComm.lean},
\texttt{HiAlgUnit.lean}, \texttt{SphComm.lean},
\texttt{LocIdem.lean}, \texttt{LocFree.lean},
\texttt{HeckeLocal.lean}, and \texttt{InvMem.lean}, at the common
source pin. The companion provenance gives their full paths and
hashes. The only external algebra facts used as checked published
applications here are finite-module completeness and finite-algebra
local decomposition over a complete Noetherian local ring:
\url{https://stacks.math.columbia.edu/tag/00MA} and
\url{https://stacks.math.columbia.edu/tag/0325}.
They are applied to rings and modules whose Noetherianity,
completeness and finiteness were verified explicitly in OHA5.
The block identification, freeness, residue field, unit and
original operator comparisons were proved above.
